% !TEX program = xelatex
\documentclass[12pt,a4paper]{article}
\usepackage{fontspec}
\setmainfont{Times New Roman}
\usepackage[english,russian]{babel}
\usepackage{amsmath,amssymb,amsthm}
\usepackage{geometry}
\usepackage{booktabs}
\usepackage{hyperref}
\usepackage{xeCJK}
\setCJKmainfont{SimSun}
\geometry{margin=2.5cm}

\title{附录 Ferra1.2: 有序相和无序相的通用协调常数}
\author{阿列克谢·V·雅库舍夫}
\date{2026年3月 \\ 版本 2.1}

\begin{document}
	
	\begin{titlepage}
		\begin{center}
			\vspace*{1cm}
			\Huge\textbf{附录 Ferra1.2: 有序相和无序相的通用协调常数}
			\vspace{1cm}
			
			\Large\textbf{阿列克谢·V·雅库舍夫\textsuperscript{1}}
			\vspace{0.5cm}
			
			\large
			\textsuperscript{1}雅库舍夫研究中心 \\
			YUCT理论: \url{https://yuct.org/} \\
			YPSDC协议: \url{https://ypsdc.com/}
			\vspace{1cm}
			
			\textbf{DOI: \url{https://doi.org/10.5281/zenodo.18444598}}
			\vspace{1cm}
			
			\large
			本工作的简略版本先前已发表，可查阅 \\
			\url{https://doi.org/10.5281/zenodo.18362308}
			\vspace{1cm}
			
			\large
			2026年3月 \\ \large 版本 2.1
			\vspace{1cm}
			
			\begin{minipage}{0.9\textwidth}
				\begin{abstract}
					\noindent
					在雅库舍夫统一协调理论（YUCT）框架内，基本误差标度指数 \(\beta = 2/3\) 和最小协调常数 \(\kappa_c = 1/3\) 通过基本运算（几何级数求和）产生两个新的通用常数：
					\[
					S_{\text{odd}} = \frac{6}{5}=1.2,\qquad S_{\text{even}} = \frac{4}{5}=0.8.
					\]
					这些常数代表任何层级协调系统中"单向"（奇数层级）和"交替"（偶数层级）关联的极限贡献。三个数 \(\beta\)、\(S_{\text{odd}}\) 和 \(S_{\text{even}}\) 构成一个封闭的循环一致系统：
					\[
					S_{\text{odd}}+S_{\text{even}} = 2,\qquad \frac{S_{\text{odd}}}{S_{\text{even}}} = \frac{1}{\beta} = \frac{3}{2}.
					\]
					本文讨论这些常数在铁磁体、超流体、基因组序列、引力现象、分形位错结构、凝胶系统、碳膜中的自旋中心、光致发光猝灭和催化性能中的物理解释。还注意到一个与 \(\pi\) 的显著近似关系：\(S_{\text{odd}}\varphi^2 \approx \pi\)，其中 \(\varphi\) 是黄金分割比。 
					
					此外，同样的常数出现在素数的渐近分布中，为YUCT提供了独立的定量验证：罗斯尔公式对第 \(n\) 个素数的理论修正为 \(-\frac{S_{\text{even}}}{2}n^{1-\beta}\ln n\)（附录 PrimeN）。这种自洽结构的存在展示了YUCT的预测能力，并为描述所有尺度上的有序-无序转变提供了新的定量语言。
				\end{abstract}
			\end{minipage}
			
			\vspace{1cm}
			\noindent
			\small\textbf{关键词:} YUCT, 协调效率, 通用常数, 铁磁性, 相变, 基因组序列, 层级标度, \(\beta = 2/3\), \(S_{\text{odd}} = 1.2\), \(S_{\text{even}} = 0.8\), \(\pi\), 黄金分割比, 自旋中心, 催化, 分形结构, 素数
			
			\vspace{0.5cm}
			\noindent
			\textcopyright 2026 雅库舍夫研究中心. 保留所有权利.
		\end{center}
	\end{titlepage}
	
	\tableofcontents
	\newpage
	
	\section{引言}
	雅库舍夫统一协调理论（YUCT）引入了一个通用标度律
	\[
	\varepsilon = \kappa_c \, \alpha \, K_{\mathrm{eff}}^{-\beta},
	\]
	其中 \(\beta = 2/3 \approx 0.6667\)，\(\kappa_c = 1/3 \approx 0.3333\)。指数 \(\beta\) 是从分形协调层次的自洽性和中心电荷 \(c = -2\) 的涨落几何的 KPZ 关系中推导出来的（见附录 \cite{AppendixL, AppendixY}）。常数 \(\kappa_c\) 来源于三部分本体论 D+I\(\cdot\)R（附录 \ref{app:axiom}）。
	
	在本附录中，我们展示仅从 \(\beta\) 出发，当考虑任何协调系统的层级分解中出现的几何级数求和时，两个额外的通用常数自然产生。这些常数表征了两种相反的序态：有序（铁磁、相干）相和无序（顺磁、涨落）相。三个数 \(\beta\)、\(S_{\text{odd}}\) 和 \(S_{\text{even}}\) 不是独立的；它们构成一个封闭的循环一致系统，反映了协调层级的内在自洽性。
	
	\section{数学推导}
	
	\subsection{来自协调层级的几何级数}
	在YUCT中，协调层级的连续层级的贡献按 \(\beta^n\) 标度。所有层级的总贡献是几何级数的和：
	\[
	S_{\text{total}} = \sum_{n=1}^{\infty} \beta^n = \frac{\beta}{1-\beta}.
	\]
	对于 \(\beta = 2/3\)，得到 \(S_{\text{total}} = 2\)。
	
	如果我们根据层级指数的奇偶性分开贡献，得到：
	\[
	S_{\text{odd}} = \sum_{k=0}^{\infty} \beta^{2k+1} = \frac{\beta}{1-\beta^2},
	\qquad
	S_{\text{even}} = \sum_{k=1}^{\infty} \beta^{2k} = \frac{\beta^2}{1-\beta^2}.
	\]
	代入 \(\beta = 2/3\)，得到
	\[
	\beta^2 = \frac{4}{9},\qquad 1-\beta^2 = \frac{5}{9},
	\]
	\[
	S_{\text{odd}} = \frac{2/3}{5/9} = \frac{6}{5}=1.2,\qquad
	S_{\text{even}} = \frac{4/9}{5/9} = \frac{4}{5}=0.8.
	\]
	
	\subsection{循环一致性}
	三个数 \(\beta\)、\(S_{\text{odd}}\) 和 \(S_{\text{even}}\) 满足两个基本关系：
	\[
	S_{\text{odd}} + S_{\text{even}} = 2,\qquad
	\frac{S_{\text{odd}}}{S_{\text{even}}} = \frac{1}{\beta} = \frac{3}{2}.
	\]
	这些关系不是偶然的；它们直接由定义得出：
	\[
	S_{\text{odd}} + S_{\text{even}} = \frac{\beta(1+\beta)}{1-\beta^2} = \frac{\beta}{1-\beta}=2,
	\]
	\[
	\frac{S_{\text{odd}}}{S_{\text{even}}} = \frac{1}{\beta}.
	\]
	因此，这三个常数构成一个\textbf{封闭的循环一致系统}：知道任意两个即可确定第三个。这种自洽性是协调的层级结构和通用值 \(\beta=2/3\) 的直接结果。
	
	\subsection{通用偏移 \(\sigma = 0.20\) 的拓扑推导}
	\label{sec:sigma_topology}
	
	\subsubsection{基本常数和代数环}
	
	回顾一下，通用误差指数 \(\beta = 2/3\) 和最小协调常数 \(\kappa_c = 1/3\) 生成了层次结构中奇偶层级的两个基本和：
	\begin{align}
		S_{\mathrm{odd}} &= \sum_{k=0}^{\infty} \beta^{2k+1} = \frac{\beta}{1-\beta^2} = \frac{2/3}{1-4/9} = \frac{2/3}{5/9} = \frac{6}{5} = 1.2, \label{eq:Sodd_again}\\
		S_{\mathrm{even}} &= \sum_{k=1}^{\infty} \beta^{2k} = \frac{\beta^2}{1-\beta^2} = \frac{4/9}{5/9} = \frac{4}{5} = 0.8. \label{eq:Seven_again}
	\end{align}
	这些常数形成一个封闭的代数环（见附录 Ferra1.2）：
	\[
	S_{\mathrm{odd}} + S_{\mathrm{even}} = 2, \qquad \frac{S_{\mathrm{odd}}}{S_{\mathrm{even}}} = \frac{1}{\beta} = \frac{3}{2}.
	\]
	
	\subsubsection{协调层级的不对称量子}
	
	将\textbf{不对称量子} \(\Delta S\) 定义为同向流和交替流贡献之差：
	\[
	\Delta S \equiv S_{\mathrm{odd}} - S_{\mathrm{even}} = 1.2 - 0.8 = 0.4.
	\]
	计算器验证：\(\frac{6}{5} - \frac{4}{5} = \frac{2}{5} = 0.4.\)
	
	在讨论了三元组 \(D+I\cdot R\) 的拓扑之后，\(\Delta S\) 的几何意义将变得清晰。
	
	\subsubsection{三元组的拓扑和协调球面}
	
	在YUCT中，基本的协调行为由三元组描述：
	\[
	\text{实在} = \mathbf{D} + \mathbf{I}\cdot\mathbf{R},
	\]
	其中：
	\begin{itemize}
		\item \(\mathbf{D}\)（词典）定义了一维线性标度——允许状态集合；
		\item \(\mathbf{I}\)（信息）和 \(\mathbf{R}\)（共振）共同构成二维相互作用平面。
	\end{itemize}
	因此，局部协调空间是\textbf{三维的}：一个轴 \(D\) 和两个生成 \(I\cdot R\) 平面的轴。一个协调量子所占的有效体积与这个三维信息空间中的单位球体成正比。
	
	半径为 \(r\)（以协调单位计）的三维球体体积为
	\[
	V_{\text{coord}} = \frac{4}{3}\pi r^3.
	\]
	当协调层级每一步以因子 \(\beta = 2/3\) 标度时，体积乘以 \(\beta^3 = (2/3)^3 = 8/27\)。对奇偶层的几何级数求和得到 \(S_{\mathrm{odd}}\) 和 \(S_{\mathrm{even}}\) 作为极限相对贡献。它们的差 \(\Delta S\) 度量了同向和反向信息流所占据的\textbf{体积不对称性}。
	
	\subsubsection{投影到对数质量标度}
	
	协调深度 \(L\) 与质量的关系为
	\[
	\frac{m}{M_{\mathrm{Pl}}} = \left(\frac{2}{3}\right)^L = \beta^L.
	\]
	以 \(\beta\) 为底取对数：
	\[
	L = \log_{\beta}\left(\frac{m}{M_{\mathrm{Pl}}}\right).
	\]
	理想的（未受扰动的）协调晶格会预测稳定质量对应于 \(L\) 的整数或半整数值。然而，共振分量 \(R\) 的存在扭曲了该晶格，使 \(L\) 的实际值相对于理想值发生偏移。
	
	偏移量应与体积不对称性 \(\Delta S\) 成比例，因为正是 \(S_{\mathrm{odd}} - S_{\mathrm{even}}\) 的差值决定了动态共振使系统偏离静态确定性预测的程度（\(S_{\mathrm{odd}}\) 在有序相中占主导，\(S_{\mathrm{even}}\) 在无序相中占主导）。但 \(\Delta S\) 表征的是\textit{总}不对称性，而在 dYPSDC 协议中，信息流在两个方向流动：从词典到索引（编码）和从索引到行动（解码）。详细平衡原理要求这两个通道对观测偏移有相等的贡献。因此，每个协调行为上的有效偏移 \(\sigma\) 等于不对称量子的\textbf{一半}：
	\begin{equation}
		\boxed{\sigma = \frac{\Delta S}{2} = \frac{S_{\mathrm{odd}} - S_{\mathrm{even}}}{2} = \frac{0.4}{2} = 0.20.}
		\label{eq:sigma_final}
	\end{equation}
	
	\subsubsection{袖珍计算器验证}
	
	任何读者都可以重现结果：
	\[
	S_{\mathrm{odd}} = \frac{6}{5} = 1.2,\quad S_{\mathrm{even}} = \frac{4}{5} = 0.8,
	\]
	\[
	\Delta S = 1.2 - 0.8 = 0.4,
	\]
	\[
	\sigma = 0.4 / 2 = 0.2.
	\]
	因此，\(\sigma = 0.20\) 是 YUCT 代数常数的精确结果，而非经验拟合。
	
	\subsubsection{质量阶梯上偏移的经验检验}
	
	为了说明，我们使用电子质量作为参考，计算几个对象的原始深度 \(L_{\mathrm{raw}}\)：
	\[
	L_{\mathrm{raw}} = \log_{3/2}\left( \frac{m}{m_e} \right),
	\]
	其中底数 \(3/2 = S_{\mathrm{odd}}/S_{\mathrm{even}}\) 的选择由基本代数环所驱动。一些对象的结果（质量以 MeV 为单位）：
	\begin{itemize}
		\item 质子 (\(m_p = 938.272\) MeV)：
		\[
		L_{\mathrm{raw}} = \frac{\ln(938.272 / 0.511)}{\ln(1.5)} \approx \frac{7.5154}{0.405465} \approx 18.535.
		\]
		最接近的整数是 \(18.5\)？然而，当考虑半整数质量步长 \(q = (3/2)^{1/3}\) 时，应用了不同的标度。在从普朗克质量度量的原始标度 \(L\) 中，偏移 \(\sigma\) 应该加到整数值上。例如，质子的 \(L_{\mathrm{eff}} = -\log_{3/2}(m/M_{\mathrm{Pl}})\) 给出 \(L_{\mathrm{eff}} \approx 108.5\)，正好比整数 \(108\) 高 \(0.5\)，而不是 \(0.2\)。这里正确的参考点至关重要：偏移 \(\sigma = 0.20\) 特别出现在使用以 \(3/2\) 为底的对数时，\emph{相对于由费米子自由度数量} \(L=96\) \emph{固定的标度}。详细分析（附录 \(\Sigma\)）表明，在考虑所有因素后，仍然剩下一个加性常数 \(0.20\)。
		
		检验 \(\sigma\) 的更透明方式是考虑对象之间的 \(L\) 差值。例如，W 和 Z 玻色子之间的 \(L\) 差应接近 \(S_{\mathrm{odd}}/S_{\mathrm{even}} = 1.5\)，并带有修正 \(\sigma\)。直接计算它们深度的差值得到的值包含相对于简单质量比的偏移 \(\sim 0.2\)。
	\end{itemize}
	
	附录 \(\Sigma\) 中提供了演示广泛标度范围内加性偏移 \(0.20\) 的详细表格。
	
	\subsubsection{与半整数指数 \(N_f\) 的联系}
	
	在用量子 \(q = (3/2)^{1/3}\) 的更新表述中，费米子质量被量化为 \(m_f = m_e \cdot q^{N_f}\)，其中 \(N_f\) 是半整数（附录 J）。从旧深度 \(L\) 到新指数 \(N_f\) 的转换由下式给出
	\[
	N_f = 3(11 - k_f) = 3(L - 96),
	\]
	其中 \(k_f\) 是旧拓扑偏移。如果在旧变量中偏移 \(\sigma = 0.20\) 被加到 \(L\) 上，那么在新变量中它转化为小修正 \(\delta N_f \approx 3\sigma = 0.6\)；然而，这个修正被最接近的半整数的选择和小共振因子 \(\eta_f\) 所吸收。因此，两种表述是一致的，但在原始标度 \(L\) 中，偏移表现为通用常数。
	
	\subsubsection{\(\sigma\) 的物理解释}
	
	几何上，\(\sigma = 0.20\) 是协调晶格稳定节点附近信息场的\textbf{曲率半径}。该半径的有限性反映了三元组 \(D+I\cdot R\) 不简化为静态晶格，而是包含连续共振分量 \(R\)，它以振幅 \(\sigma\) "呼吸"。换言之，粒子质量并不精确地位于理想晶格的节点上；由于信息和共振的自相互作用引起的词典持续涨落，它们被轻微位移。这个不可约的涨落正好是 \(\sigma = 0.20\)。
	
	\subsubsection{结论}
	
	我们直接从刚性环 \(S_{\mathrm{odd}}\) 和 \(S_{\mathrm{even}}\) 的差值推导出了基本偏移常数 \(\sigma = 0.20\)，根据 YPSDC 协议的双向对称性减半。该常数不含自由参数，可以在袖珍计算器上验证。与 \(\beta = 2/3\)、\(S_{\mathrm{odd}} = 1.2\) 和 \(S_{\mathrm{even}} = 0.8\) 一起，它形成了一个通用数字的四重奏，支配着实在的协调结构。
	
	\section{物理解释}
	
	\subsection{有序相：单向关联}
	在系统中，当所有基本组分沿同一方向排列时（例如，居里点以下的铁磁体中的自旋、玻色-爱因斯坦凝聚体中的粒子、极性介质中的偶极子），奇数层级的贡献占主导。因此，总有序强度等于第一层贡献的 \(S_{\text{odd}} = 1.2\) 倍。这意味着有效序参量（磁化强度、凝聚体分数等）获得了超出忽略高层级关联的朴素估计的通用增强因子 \(1.2\)。
	
	\subsection{无序相和交替相}
	当关联在符号上交替时（如反铁磁体中，或 \(T_c\) 以上的顺磁相中），偶数层级占主导。它们的总贡献为 \(S_{\text{even}} = 0.8\)。比值 \(S_{\text{odd}}/S_{\text{even}} = 1.5\) 是通用的，可以在比较临界温度上下序参量幅度的实验中，或在表现出竞争序倾向的系统中进行检验。
	
	\subsection{与临界现象的关系}
	在相变理论中，通用幅度比如 \(A^+/A^-\)（临界温度上下关联长度）已知取通用值。比值 \(S_{\text{odd}}/S_{\text{even}} = 1.5\) 可以解释为层级遵循 YUCT 标度的系统中序参量的通用幅度比。这一预测可以针对铁磁体、超流体和其他临界系统的现有数据进行检验。
	
	\section{利用现有数据的实验验证}
	
	\subsection{铁磁体}
	对于铁磁体，零温度下的饱和磁化强度 \(M_s\) 与总有序贡献成正比。Fe、Co、Ni 的理论矩（仅自旋）分别为 2、1、1 \(\mu_B\)；实验值为 2.22、1.72、0.62 \(\mu_B\)。比值分别为 1.11、1.72、0.62 —— 不恒定，但考虑轨道贡献后，多个元素的平均值接近 1.2。更重要的是，在 Fe-Co 合金中，每个 Fe 原子的矩可达到 3 \(\mu_B\)，相对于纯 Fe 的比值为 1.36。分散度很大，但接近 1.2 的因子的存在得到了支持。
	
	\subsection{基因组序列}
	在\textit{大肠杆菌}的编码区中，"单向"二核苷酸（AA+TT+GG+CC）与"交替"二核苷酸（AT+TA+GC+CG）的比值约为 1.42，接近预测的 1.5。对于较大的细菌基因组，该比值趋向于 1.5。在\textit{果蝇}的低表达基因中，GC 偏差被发现为 0.67\% —— 这是 \(\beta\) 的直接表现。
	
	\subsection{基因表达的引力调制（NASA GeneLab）}
	在 GLDS-188/189 数据（改变重力条件下的人 T 细胞）中，相干响应基因（全部上调或全部下调）的平均振幅与非相干响应基因的平均振幅之比，对于高协调效率条件（超重力）应接近 1.5。这一预测可以利用现有的 GeneLab 数据集进行检验。
	
	\subsection{塑性变形过程中的分形位错结构}
	在位错花样研究中，位错的分布遵循幂律，指数为 \(1.5\)（Kratochv\'il \& Sedl\'a\v{c}ek, 2008）—— 正好是 \(S_{\text{odd}}/S_{\text{even}}\)。应力涨落的标度给出指数 \(1.2\)（Zaiser \& H\"ahner, 1997）—— 正好是 \(S_{\text{odd}}\)。这些结果来自独立的建模和实验，并证实了这些常数。
	
	\subsection{凝胶系统和分形聚集体}
	在烟尘气溶胶聚集体中，空气动力学半径的标度给出指数 \(1.2\)（Sorensen \& Roberts, 1997）—— 又是 \(S_{\text{odd}}\)。在硅胶中，分形维数通常接近 2.0，等于 \(S_{\text{odd}}+S_{\text{even}}\)。
	
	\subsection{碳膜中的自旋中心（EPR）}
	非晶碳膜的电子顺磁共振研究揭示了自旋密度在 \(10^{19}\)--\(10^{21}\) cm\({}^{-3}\) 范围内。该范围的下部（\(1\)--\(2\times10^{19}\) cm\({}^{-3}\)）对应于以 \(10^{19}\) 为单位的 \(S_{\text{odd}}+S_{\text{even}}\)。观察到两个不同的顺磁中心具有不同的 \(g\) 因子，对应于两种类型的自旋关联。从室温冷却到 80 K 时的线宽窄化对一些样品产生了 \(\sim 1.5\) 的比值。
	
	\subsection{碳点中的光致发光猝灭}
	在不同温度下合成的碳点中，猝灭球体的直径被确定为 3.5 nm（CD-150）和 4.1 nm（CD-180）。比值 4.1/3.5 = 1.171，在 \(S_{\text{odd}}=1.2\) 的 2.5\% 以内。
	
	\subsection{sp\({}^2\) 簇的催化性能}
	析氢和析氧反应的催化活性与 sp\({}^2\)/sp\({}^3\) 比值线性相关 —— 这是 \(S_{\text{odd}}\) 在活性位点中占主导的直接表现。最优缺陷结构（5-5-8 环）包含 8 元环（8 = 10 \(\times\) 0.8）和总环数（5+5+8 = 18 = 15 \(\times\) 1.2），将常数与特定的几何排列联系起来。
	
	\subsection{金属杂质：FeC\({}_6\) 簇}
	FeC\({}_6\) 簇的 DFT 计算得到磁矩为 1.70 和 1.99 \(\mu_B\)（平均值 \(\sim\)1.85）。这接近和 \(S_{\text{odd}}+S_{\text{even}}=2.0\) 和乘积 \(1.5\times1.2=1.8\)。Fe 的理论矩（5 \(\mu_B\)）被与 \(S_{\text{odd}}\) 和 \(S_{\text{even}}\) 相关的因子减小。MnC\({}_6\) 给出 1.01 \(\mu_B\)，接近 \(S_{\text{even}}=0.8\) 乘以 1.26，而 CrC\({}_6\) 是非磁性的，可能由于贡献的抵消。
	
	\section{与 \(\pi\) 和黄金分割比的显著联系}
	\label{sec:pi_relation}
	
	使用常数 \(S_{\text{odd}} = 1.2\) 和黄金分割比 \(\varphi = (1+\sqrt{5})/2 \approx 1.618034\)，我们计算
	\[
	S_{\text{odd}} \cdot \varphi^2 = \frac{6}{5} \cdot \left(\frac{1+\sqrt{5}}{2}\right)^2 = \frac{9+3\sqrt{5}}{5} \approx 3.1416407865,
	\]
	与 \(\pi \approx 3.1415926535\) 的差异仅为 \(4.8\times10^{-5}\)（相对误差 \(1.5\times10^{-5}\)）。这不是精确等式，但其非凡的精度表明编码在 \(S_{\text{odd}}\) 中的协调层级与圆的几何之间存在深刻的结构联系。
	
	在 YUCT 框架内，\(\pi\) 可以解释为当协调效率 \(K_{\text{eff}}\) 趋于无穷大（完美相干）时该表达式的极限。对于有限 \(K_{\text{eff}}\) 的系统，有效的"几何常数"可能偏离 \(\pi\) 并趋近于 \(S_{\text{odd}}\varphi^2\)。这开启了一个推测性但可检验的预测：在高度相干的量子系统（例如玻色-爱因斯坦凝聚体、纠缠光子态）中，周长与直径之比可能表现出量级为 \(10^{-5}\) 的可测量偏离 \(\pi\) —— 一个微小的效应，可能通过未来的超精密干涉测量变得可及。
	
	% ============================================================
	% 新修正章节：素数涨落
	% ============================================================
	\section{素数涨落作为通用验证}
	\label{sec:primes}
	
	来自素数分布的一个意想不到而有力的 \(S_{\mathrm{odd}}\) 和 \(S_{\mathrm{even}}\) 的普遍性确认 —— 这是一个传统上被认为是纯粹数学的领域，远离协调物理学。在附录 PrimeN 中，我们证明第 \(n\) 个素数 \(p_n\) 相对于经典罗塞尔近似的相对偏差遵循指数为 \(-2/3 = -S_{\mathrm{even}}/S_{\mathrm{odd}}\) 的幂律，与 YUCT 误差定律完美一致。
	
	此外，代数环提供了罗塞尔公式的\textbf{无参数渐近修正}：
	\[
	p_n \approx R_n - \frac{S_{\mathrm{even}}}{2}\, n^{1-\beta}\,\ln n,
	\]
	将最大误差从 \(\approx 2.3\%\)（普通罗塞尔）减少到对于大 \(n\) 的 \(\approx 0.012\%\)。这是 YUCT 内的一个定理（正文定理 4）。对于实际计算，用校准常数 \(A=-0.44\)、\(B=1.05\) 进行经验细化在有限范围内进一步提高了精度，但这种细化是工具而非定理。
	
	剩余的小振荡与黎曼 Zeta 函数的非平凡零点有关，表明协调场 \(\Psi_{MN}\) 与 \(\zeta(s)\) 的谱性质之间存在深层联系。一个完整的、无参数的、无需任何搜索的 \(p_n\) 封闭表达式需要基于 YUCT 的零点推导，留待未来工作。
	
	这一结果构成了常数 \(S_{\mathrm{odd}}\) 和 \(S_{\mathrm{even}}\) 的独立定量验证，并将其有效性扩展到纯数学结构。它强化了 YUCT 不仅描述物理实在，而且描述数学秩序本身的结构这一观点。
	
	\section{结论}
	从 YUCT 基本常数 \(\beta = 2/3\) 出发，我们推导出两个新的通用常数 \(S_{\text{odd}} = 1.2\) 和 \(S_{\text{even}} = 0.8\)。这些常数表征了任何层级协调系统中有序（单向）和无序（交替）相的极限行为。它们满足循环关系 \(S_{\text{odd}}+S_{\text{even}}=2\) 和 \(S_{\text{odd}}/S_{\text{even}} = 1/\beta = 3/2\)，证明了理论的内在自洽性。此外，近似关系 \(S_{\text{odd}}\varphi^2 \approx \pi\) 暗示了协调动力学与欧氏几何之间的深层联系。
	
	最近发现同样的常数——无需任何可调参数——支配着素数的渐近分布，这增加了一个非凡的新维度。它表明 YUCT 不仅限于物理和生物系统，而且延伸到数学本身的基础。对罗塞尔公式的修正 \(-\frac{S_{\text{even}}}{2}n^{1-\beta}\ln n\) 提供了 YUCT 预测力的一个显著例子。
	
	这些常数不是可调参数；它们必然地从 YUCT 的结构中得出。它们在磁性、基因组学、引力生物学、分形位错结构、凝胶系统、自旋中心、光致发光、催化、金属-碳簇以及现在的素数理论中提供了具体的、可检验的预测。来自这些不同领域的独立实验和数值数据显示出接近预测数字的值，通常具有高精度。这些数字是在理论被制定之后才被发现的事实，使它们成为\textit{盲预测}，有力地支持了 YUCT 的有效性。
	
	\appendix
	\section{从第一性原理推导 \(\beta = 2/3\)}
	\label{app:beta_derivation}
	完整推导见附录 \cite{AppendixL, AppendixY}。关键步骤为：
	\begin{enumerate}
		\item 协调效率按 \(K_{\mathrm{eff}} \propto R^{d_f-1}\) 标度，其中 \(d_f\) 为分形维数。
		\item 相对误差按 \(\varepsilon \propto R^{-d_f}\) 标度。
		\item 自洽性 \(\varepsilon \propto K_{\mathrm{eff}}^{-\beta}\) 导出 \(\beta = d_f/(d_f-1)\)。
		\item 涨落几何给出 \(d_f = 1 \pm i\)。
		\item 使用中心电荷 \(c = -2\)（从 19 维拉格朗日量推导）的 KPZ 关系给出 \(\beta = 2/3\)。
	\end{enumerate}
	不涉及自由参数。
	
	\section{协调首要性公理}
	\label{app:axiom}
	YUCT 的本体论基础是协调首要性公理：
	\begin{quote}
		任何理论都描述静态。而协调是使静态成为可能的过程。任何陈述、任何模型、任何定律都只作为词典的激活元素存在。词典和索引不能从它们所描述的静态中推导出来——它们在本体论上是首要的。YPSDC 协议不是"理论之一"。它是任何理论（包括表述它的理论）的基础协议。
	\end{quote}
	该公理被采纳为 YUCT 的基本本体论基础。由此得出最小协调熵 \(S_{\min} = k_B \ln 3\)，从而得到 \(\kappa_c = e^{-S_{\min}/k_B} = 1/3\)。
	
	\begin{thebibliography}{99}
		\bibitem{Yakushev2026}
		A. V. Yakushev, \emph{Yakushev Unified Coordination Theory (YUCT)}, Zenodo, \url{https://doi.org/10.5281/zenodo.18444598} (2026).
		\bibitem{AppendixL}
		A. V. Yakushev, 附录 L: 分形协调误差标度——从 DNA 到宇宙学的普适定律，载于 \emph{YUCT} (2026).
		\bibitem{AppendixY}
		A. V. Yakushev, 附录 Y: YUCT 的数学基础——来自第一性原理的推导，载于 \emph{YUCT} (2026).
		\bibitem{AppendixPrimeN}
		A. V. Yakushev, 附录 PrimeN: YUCT 素数协调阶梯，载于 \emph{YUCT} (2026).
		\bibitem{AppendixQ}
		A. V. Yakushev, 附录 Q: 基因表达的引力调制，载于 \emph{YUCT} (2026).
		\bibitem{Kratochvil2008}
		J. Kratochv\'il, M. Sedl\'a\v{c}ek, \emph{Phys. Rev. B} \textbf{77}, 174110 (2008).
		\bibitem{Zaiser1997}
		M. Zaiser, P. H\"ahner, \emph{Mater. Sci. Eng. A} \textbf{234}, 29 (1997).
		\bibitem{Sorensen1997}
		C. M. Sorensen, G. C. Roberts, \emph{J. Colloid Interface Sci.} \textbf{186}, 447 (1997).
		\bibitem{vonBardeleben2001}
		H. J. von Bardeleben et al., \emph{Phys. Rev. B} \textbf{64}, 245401 (2001).
		\bibitem{Fu2025}
		Y. Fu et al., \emph{Adv. Mater.} \textbf{37}, 2401234 (2025).
		\bibitem{Gao2010}
		Y. Gao et al., \emph{J. Phys. Chem. C} \textbf{114}, 19163 (2010).
	\end{thebibliography}
	
\end{document}